\capitulo{2}{Objetivos del proyecto}

\section{Objetivo general}

El objetivo general es diseñar e implementar una plataforma web que permita localizar aplicaciones, seleccionar uno o varios instaladores y descargarlos agrupados desde sus fuentes oficiales, evitando el almacenamiento permanente de ejecutables y ofreciendo información clara sobre el origen y el estado de cada archivo.

\section{Objetivos funcionales}

\begin{enumerate}
	\item Construir un catálogo administrable de software con nombre, descripción, etiquetas, página oficial, sistemas operativos, arquitecturas y estado de disponibilidad.
	\item Permitir búsquedas parciales por nombre y filtros por etiquetas. El filtrado podrá exigir todas las etiquetas seleccionadas o un mínimo configurado por el usuario.
	\item Permitir la selección de una o varias aplicaciones. Una selección múltiple se tratará como un bundle temporal para generar un ZIP.
	\item Crear bundles oficiales, públicos y privados. Los usuarios anónimos podrán personalizar temporalmente los bundles públicos; los usuarios registrados podrán, además, guardar y publicar sus propios bundles.
	\item Asignar un enlace estable a cada bundle público y ordenar el catálogo de bundles por popularidad, fecha, nombre o etiquetas.
	\item Incorporar una estrella binaria por usuario y bundle. Cada usuario registrado podrá añadirla o retirarla una vez, y el número total servirá para ordenar los bundles.
	\item Permitir a los usuarios registrados solicitar la incorporación de una aplicación, proporcionando al menos su nombre y su página oficial.
	\item Proporcionar al administrador funciones para mantener aplicaciones, etiquetas, reglas de resolución, dominios permitidos y estados del catálogo.
	\item Informar del progreso de cada descarga, de los archivos incluidos y de los errores que impidan completar una parte de la solicitud.
\end{enumerate}

\section{Objetivos de descarga y procesamiento}

\begin{enumerate}
	\item Resolver dinámicamente el enlace del instalador desde una fuente oficial mediante una estrategia específica o, como respaldo, mediante navegación automatizada con Playwright.
	\item Validar protocolo, dominio final, redirecciones, tipo de contenido, extensión y tamaño antes de aceptar un archivo.
	\item Descargar los instaladores en el servidor, generar un ZIP con un manifiesto y eliminar todos los archivos temporales al terminar, cancelar o expirar el trabajo.
	\item Procesar las solicitudes de manera asíncrona con RabbitMQ para separar la interacción web de las operaciones de resolución, descarga y compresión.
	\item Aplicar reintentos controlados y permitir un resultado parcial cuando una aplicación falle, manteniendo una explicación individual para cada incidencia.
	\item Calcular, cuando las fuentes lo permitan, un tamaño estimado antes de iniciar el trabajo y aplicar límites por archivo, ZIP, usuario, dirección IP y número de operaciones concurrentes.
\end{enumerate}

\section{Objetivos de búsqueda semántica}

\begin{enumerate}
	\item Representar las aplicaciones mediante embeddings generados a partir de sus metadatos y almacenarlos en PostgreSQL mediante pgvector.
	\item Interpretar consultas en lenguaje natural y recuperar candidatos semánticamente próximos antes de aplicar los filtros ordinarios del catálogo.
	\item Utilizar un modelo de lenguaje para proponer bundles, explicar recomendaciones y resumir errores a partir de información recuperada del sistema.
	\item Validar las salidas del modelo mediante esquemas estrictos y mantener desacoplado el proveedor para poder sustituir Groq por un modelo local o por otro servicio.
	\item Impedir que el modelo resuelva o apruebe enlaces, cambie dominios autorizados, ejecute descargas o modifique reglas sin revisión administrativa.
\end{enumerate}

\section{Objetivos técnicos y de calidad}

\begin{itemize}
	\item Desarrollar un cliente ligero con Vite, HTML, CSS y TypeScript.
	\item Construir el backend con Java 26 y Spring Boot, separando las responsabilidades en servicios desplegables y escalables.
	\item Documentar la API HTTP mediante OpenAPI, cuyo propósito es describir de forma independiente del lenguaje las capacidades de una API~\cite{openapi}.
	\item Utilizar MySQL para la información transaccional y PostgreSQL con pgvector para la búsqueda vectorial.
	\item Contenerizar los componentes con Docker y preparar su despliegue compatible con Compose.
	\item Definir pruebas unitarias, de integración, de seguridad, de resolución de enlaces, de generación de ZIP y de extremo a extremo.
	\item Mantener trazabilidad mediante registros de administración, resolución, descargas y uso del asistente.
\end{itemize}

\section{Objetivos de seguridad y sostenibilidad}

El sistema deberá reducir los riesgos derivados de procesar URL externas y ejecutables. Para ello se definirán listas de dominios oficiales, bloqueo de redes privadas, límites de redirección, sanitización de nombres, protección frente a SSRF, control de acceso y cuotas de uso. La plataforma no ejecutará los instaladores descargados.

Desde el punto de vista de sostenibilidad, se buscará evitar copias permanentes, eliminar temporales, limitar trabajos innecesarios, reutilizar metadatos válidos y escalar los servicios según su carga. También se atenderá a la accesibilidad, la claridad de la información, la protección de datos y el acceso equitativo a la funcionalidad básica.
